\documentclass[10pt,a4paper]{article}
\usepackage{amsmath,amssymb,amsthm,graphicx}
\title{Dense transcendental power series with easy truncations and a factorial-hard comparison}
\author{Pranjal Dutta, Gorav Jindal, Anurag Pandey, Amit Sinhababu}
\date{doi:10.4230/LIPIcs.CCC.2021.25; CC BY4.0}
\begin{document}
\maketitle
\noindent\textit{Editorial note.} Original authors: Pranjal Dutta, Gorav Jindal, Anurag Pandey, Amit Sinhababu. Official CCC2021 article25, DOI10.4230/LIPIcs.CCC.2021.25.\par\medskip
\noindent\textit{Editorial note.} License: https://creativecommons.org/licenses/by/4.0/\par\medskip
\noindent\textit{Editorial note.} Faithful editable TeX transcription of formal author Theorems38/40/44 with complete Section7 proofs, exact sequence/complexity definitions and cited transcendence/functional-equation premises, and full published bibliography. The direct factorial contrast is grouped with the two dense constructions. Original author expressions and conditions preserved without mathematical correction or new proofs; only official CC BY4.0 version used.\par\medskip
\noindent\textit{Allowed standard prerequisites (editorial).}\begin{itemize}
\item Stern generating function identity and transcendence. Exact author Theorem37[12] retained.
\item Transcendence of the multiplicative series. Exact author Theorem39[13] retained.
\item Zero-radius criterion for formal algebraic series. Exact author Lemma42[28] retained.
\end{itemize}\medskip
\section*{Author power-series and complexity conventions}
A formal power series is a generalization of a polynomial, where the number of terms can be infinite. Formally, $A=\sum_{i\ge0}A_ix^i$ with $A_i\in\mathbb F$, is a power series in the power series ring $\mathbb F[[x]]$. We define the degree $d$ truncation $\mathrm{trunc}(A,d):=\sum_{0\le i\le d}A_ix^i$.
\paragraph{Definition 12 (Algebraic and Transcendental Power Series).} A formal power series $f\in\mathbb C[[x]]$ is said to be algebraic if there exists a polynomial $g\in\mathbb C[x][t]$ such that $g(f)=0$. Otherwise $f$ is said to be transcendental.
For a polynomial $f\in\mathbb F[\mathbf x]$, $L_{\mathbb F}(f)$ is the length of the smallest division-free arithmetic circuit (with only $\{+,-,\times\}$ gates) computing $f$. Constant-free complexity $\tau(f)$ is the size of the minimal constant-free circuit computing $f$, with constants in $\{-1,0,1\}$. Straight-line programs for integers start at 1 and use $+,-,\times$; $\tau(N)$ is the minimum such length for $N$. Subscripts are omitted when the field is clear. These are the author conventions from Sections2/5.2.
\paragraph{Definition 36 (The Stern sequence).} The sequence $(a_n)_{n\ge0}$ given by $a_0=0,a_1=1$, and when $n\ge1$, by $a_{2n}=a_n$ and $a_{2n+1}=a_n+a_{n+1}$, is called the Stern sequence. Its generating function is $A(x):=\sum a_nx^n$.
\paragraph{Theorem 37 (Lemma2.1 and Theorem2.2 in [12]).} If $A(z)$ is the generating function of the Stern sequence, then (1) $A(x^2)=A(x)\frac{x}{x^2+x+1}$; (2) the function $A(x)$ is transcendental.
\paragraph{Multiplicative sequence.} The sequence $(f_n)_{n\ge0}$ is defined as: $f_0=1,f_1=1,f_2=-1,f_p=1$ for all odd primes $p$ and $f_{ab}=f_af_b$. We look at the corresponding generating function $F(x):=\sum f_nx^n$.
\paragraph{Theorem 39 ([13,Theorem2]).} The power series $F(x)$ is transcendental.
\paragraph{Holonomic sequence.} A sequence $(h_n)_{n\ge0}$ is called holonomic if it satisfies a recurrence of the form $a_r(n)h_{n+r}+a_{r-1}(n)h_{n+r-1}+\cdots+a_0(n)h_n=0$, where $a_i$ are polynomials in $n$. The corresponding generating function $H(x):=\sum h_nx^n$ is said to be a holonomic function.
\paragraph{Lemma 42 ([28]).} If $F(x)=\sum_{n\ge0}f_nx^n$ is a power series in $\mathbb C[[x]]$ and the radius of convergence of $F(x)$ is zero then $F(x)$ is transcendental. This follows directly from Proposition2 in [28], as stated by the authors.

\section*{Formal author transcendental-truncation examples}
\paragraph{Theorem 38.} For the generating function $A(x)$ of the Stern sequence, we have $L(\mathrm{trunc}(A(x),d))=O(\log^2d)$.
\paragraph{Theorem 40.} For $F(x)$ from the multiplicative sequence in Section7.1.2, we have $L(\mathrm{trunc}(F(x),d))=O(\log^2d)$.
\paragraph{Theorem 44.} For $F(x)=\sum_{n\ge0}n!x^n$, if $\tau(\mathrm{trunc}(F(x),d))=\mathrm{poly}(\log d)$ then $(d!)$ has complexity $\mathrm{poly}(\log d)$.

\section*{Complete author Section 7.1.1 proof}
\paragraph{Proof of Theorem 38.} By using Theorem 37, we obtain that:
\[A(x)=(x^2+1)A(x^2)+\frac{A(x^2)}x.\tag{3}\]
Suppose $B_d(x):=\mathrm{trunc}(A(x),\lfloor d/2\rfloor+1)$. Notice that the degree of $C_d(x):=(x^2+1)B_d(x^2)+B_d(x^2)/x$ is at most $2\lfloor d/2\rfloor+4$ and $\mathrm{trunc}(C_d(x),d)=\mathrm{trunc}(A(x),d)$. Hence we can compute $\mathrm{trunc}(A(x),d)$ from $C_d(x)$ by subtracting at most 4 monomials, which can be done using $O(\log d)$ gates. Also $B_d(x)$ can be computed from $\mathrm{trunc}(A(x),\lfloor d/2\rfloor)$ using $O(\log d)$ gates. Hence we obtain the following recurrence:
\[L(\mathrm{trunc}(A(x),d))\le L(\mathrm{trunc}(A(x),\lfloor d/2\rfloor))+O(\log d).\]
This implies, $L(\mathrm{trunc}(A(x),d))=O(\log^2d)$.\qed
\section*{Complete author Section 7.1.2 proof}
\paragraph{Proof of Theorem 40.} We use the notation $\nu_2(m)$ to denote the highest power of 2 which divides $m\in\mathbb N$. We partition the set $[d]$ into $\lfloor\log d\rfloor$ sets $S_0,S_1,S_2,\ldots,S_{\lfloor\log d\rfloor}$ such that $k\in S_i$ iff $\nu_2(k)=i$. We define the set $O_m:=\{k:k\le m\text{ and }k\text{ is odd}\}$. Now, notice that $S_i=\{2^ik:k\in O_{\lfloor d/2^i\rfloor}\}$. For a set $S\in\mathbb N$, we define the polynomial $g_S:=\sum_{i\in S}x^i$. Observe that:
\[\mathrm{trunc}(F(x),d)=1+\sum_{i=1}^{\lfloor\log d\rfloor}(-1)^ig_{S_i}.\]
Trivially, $g_{S_i}=g_{O_{\lfloor d/2^i\rfloor}}(x^{2^i})$. Also notice that $g_{O_m}=g_{[m]}-g_{\lfloor m/2\rfloor}(x^2)$. Therefore, $L(g_{O_m})=(\log m)$, which implies that $g_{S_i}=O(\log d)$. Hence, $L(\mathrm{trunc}(F(x),d))=O(\log^2d)$.\qed
\paragraph{Remark 41.} Note that, there are power series like $\sum_{i\ge0}x^{i!}$ which are transcendental and their truncations up to degree $d$ are easy to compute. However, the series is highly sparse and degree-$d$ truncations has only $\mathrm{poly}(\log d)$ monomials, hence the easiness is trivial. The examples we discover in this work are of dense power series.
\section*{Complete author Section 7.2 comparison}
Consider the holonomic sequence $f_n=(n!)$ defined by $f_0=1$ and $f_{n+1}-(n+1)f_n=0$. Also consider the corresponding generating function $F(x)=\sum_{n\ge0}n!x^n$. We now show that $F(x)$ is transcendental and that truncation of $F(x)$ is (conditionally) hard to compute.
\paragraph{Corollary 43.} The power series $F(x)=\sum_{n\ge0}n!x^n$ is transcendental.
\begin{proof}It is clear that the radius of convergence of $F(x)$ is zero (follows from the ratio test). Hence Lemma 42 implies that $F(x)$ is transcendental.\end{proof}
\paragraph{Proof of Theorem 44.} We know that $d!x^d=\mathrm{trunc}(F(x),d)-\mathrm{trunc}(F(x),d-1)$. Setting $x=1$, we conclude.\qed

\begin{thebibliography}{99}
\bibitem{ref1} Alexander Alder. Grenzrang und Grenzkomplexität aus algebraischer und topologischer Sicht. PhD thesis, Zentralstelle der Studentenschaft, 1984.
\bibitem{ref2} Eric Allender, Peter Bürgisser, Johan Kjeldgaard-Pedersen, and Peter Miltersen. On the Complexity of Numerical Analysis. SIAM Journal on Computing, 38, January 2006. Preliminary version in the 21st Annual IEEE Conference on Computational Complexity (CCC’06). doi: 10.1109/CCC.2006.30.
\bibitem{ref3} Robert Andrews. Algebraic Hardness Versus Randomness in Low Characteristic. In 35th Computational Complexity Conference (CCC 2020), volume 169 of Leibniz International Proceedings in Informatics (LIPIcs), pages 37:1–37:32. Schloss Dagstuhl–Leibniz-Zentrum für Informatik, 2020. doi:10.4230/LIPIcs.CCC.2020.37.
\bibitem{ref4} Jean Berstel and Christophe Reutenauer. Rational Series and Their Languages. Springer- Verlag, Berlin, Heidelberg, 1988. URL: https://dl.acm.org/doi/book/10.5555/52107.
\bibitem{ref5} Markus Bläser and Gorav Jindal. On the Complexity of Symmetric Polynomials. In 10th Innovations in Theoretical Computer Science Conference (ITCS’19), volume 124 of LIPIcs, pages 47:1–47:14. Schloss Dagstuhl - Leibniz-Zentrum für Informatik, 2019. doi:10.4230/ LIPIcs.ITCS.2019.47.
\bibitem{ref6} Alin Bostan, Gilles Christol, and Philippe Dumas. Fast computation of the Nth term of an algebraic series over a finite prime field. In Proceedings of the ACM on International Symposium on Symbolic and Algebraic Computation (ISSAC’16), pages 119–126, 2016. doi: 10.1145/2930889.2930904.
\bibitem{ref7} Karl Bringmann, Christian Ikenmeyer, and Jeroen Zuiddam. On Algebraic Branching Programs of Small Width. J. ACM, 65(5):1–29, 2018. (Preliminary version in the 32nd Computational Complexity Conference (CCC’17). doi:10.1145/3209663.
\bibitem{ref8} Peter Bürgisser. The complexity of factors of multivariate polynomials. Foundations of Computational Mathematics, 4(4):369–396, 2004. arXiv:1812.06828.
\bibitem{ref9} Peter Bürgisser. On defining integers and proving arithmetic circuit lower bounds. Computa- tional Complexity, 18(1):81–103, 2009. doi:10.1007/s00037-009-0260-x.
\bibitem{ref10} Peter Bürgisser, Michael Clausen, and Amin Shokrollahi. Algebraic complexity theory, volume 315. Springer Science \& Business Media, 2013.
\bibitem{ref11} David V Chudnovsky and Gregory V Chudnovsky. On expansion of algebraic functions in power and Puiseux series, I. Journal of Complexity, 2(4):271–294, 1986. URL: https: //www.sciencedirect.com/science/article/pii/0885064X86900063.
\bibitem{ref12} Michael Coons. The Transcendence of Series Related to Stern’s Diatomic Sequence. Interna- tional Journal of Number Theory, 06, November 2011. doi:10.1142/S1793042110002958.
\bibitem{ref13} Michael Coons and Peter Borwein. Transcendence of power series for some number theoretic functions. Proceedings of the American Mathematical Society, 137, July 2008. doi:10.1090/ S0002-9939-08-09737-2.
\bibitem{ref14} Wellington De Melo and Benar Fux Svaiter. The cost of computing integers. Proceedings- American Mathematical Society, 124:1377–1378, 1996. URL: https://www.ams.org/journals/ proc/1996-124-05/S0002-9939-96-03173-5/S0002-9939-96-03173-5.pdf.
\bibitem{ref15} Richard A. Demillo and Richard J. Lipton. A probabilistic remark on algebraic program testing. Information Processing Letters, 7(4):193–195, 1978. URL: https://www.sciencedirect.com/ science/article/abs/pii/0020019078900674.
\bibitem{ref16} Pranjal Dutta. Real tau-Conjecture for sum-of-squares: A unified approach to lower bound and derandomization. In 16th International Computer Science Symposium in Russia (CSR 2021), 2021. URL: https://drive.google.com/file/d/1X8eo9GM4SCNsC2vWjPbUwMX0vff5i2k3/ view.
\bibitem{ref17} Pranjal Dutta, Nitin Saxena, and Amit Sinhababu. Discovering the roots: Uniform closure results for algebraic classes under factoring. In Proceedings of the 50th Annual ACM SIGACT Symposium on Theory of Computing, pages 1152–1165, 2018. URL: https://www.cse.iitk. ac.in/users/nitin/papers/factor-closure.pdf.
\bibitem{ref18} Pranjal Dutta, Nitin Saxena, and Thomas Thierauf. A Largish Sum-Of-Squares Implies Circuit Hardness and Derandomization. In 12th Innovations in Theoretical Computer Science Conference (ITCS 2021), volume 185 of Leibniz International Proceedings in Informatics (LIPIcs), pages 23:1–23:21. Schloss Dagstuhl–Leibniz-Zentrum für Informatik, 2021. URL: 10.4230/LIPIcs.ITCS.2021.23.
\bibitem{ref19} Joshua A Grochow et al. Complexity in ideals of polynomials: questions on algebraic complexity of circuits and proofs. Bulletin of EATCS, 2(130), 2020. URL: http://bulletin.eatcs.org/ index.php/beatcs/article/view/607.
\bibitem{ref20} Joshua A. Grochow, Ketan D. Mulmuley, and Youming Qiao. Boundaries of VP and VNP. In 43rd International Colloquium on Automata, Languages, and Programming (ICALP 2016), volume 55, pages 34:1–34:14, 2016. URL: https://core.ac.uk/download/pdf/62922137.pdf.
\bibitem{ref21} Gorav Jindal. On approximate polynomial identity testing and real root finding. PhD thesis, Saarland University, 2019. doi:10.22028/D291-29880.
\bibitem{ref22} Erich Kaltofen. Uniform closure properties of p-computable functions. In Proceedings of the eighteenth annual ACM symposium on Theory of computing, pages 330–337, 1986. doi: 10.1145/12130.12163.
\bibitem{ref23} Erich Kaltofen. Single-factor Hensel lifting and its application to the straight-line complexity of certain polynomials. In Proceedings of the 19th annual ACM symposium on Theory of computing (STOC’87), pages 443–452, 1987. doi:10.1145/28395.28443.
\bibitem{ref24} Pascal Koiran. Valiant’s model and the cost of computing integers. computational complexity, 13(3):131–146, 2005.
\bibitem{ref25} Pascal Koiran. Shallow circuits with high-powered inputs. Innovations in Computer Sci- ence (ICS), 2011. URL: https://hal-ens-lyon.archives-ouvertes.fr/ensl-00477023v4/ document.
\bibitem{ref26} Pascal Koiran and Sylvain Perifel. Interpolation in Valiant’s theory. Computational Complexity, 20(1):1–20, 2011. doi:10.1007/s00037-011-0002-8.
\bibitem{ref27} Pascal Koiran, Natacha Portier, Sébastien Tavenas, and Stéphan Thomassé. A $\tau$ -Conjecture for Newton Polygons. Foundations of computational mathematics, 15(1):185–197, 2015. doi: 10.1007/s10208-014-9216-x.
\bibitem{ref28} FV Kuhlmann. On convergent power series, 1996. URL: https://www.mathi.uni-heidelberg. de/~roquette/KONVPOTREIHEN.pdf.
\bibitem{ref29} Hsiang Kung and Joseph Traub. All Algebraic Functions Can Be Computed Fast. J. ACM,
\bibitem{ref30} Dick Lipton and Ken Regan. Factoring and factorials, February 2009. URL: https://rjlipton. wordpress.com/2009/02/23/factoring-and-factorials/.
\bibitem{ref31} Richard J Lipton. Polynomials with 0-1 coefficients that are hard to evaluate. SIAM Journal on Computing, 7(1):61–69, 1978. Preliminary version in the 16th Annual Symposium on Foundations of Computer Science (FOCS 1975). URL: https://epubs.siam.org/doi/abs/ 10.1137/0207004?journalCode=smjcat.
\bibitem{ref32} Richard J Lipton. Straight-line complexity and integer factorization. In International Algorithmic Number Theory Symposium (ANTS 94), pages 71–79. Springer, 1994. doi: 10.1007/3-540-58691-1\_45.
\bibitem{ref33} Richard J Lipton and Larry J Stockmeyer. Evaluation of polynomials with super- preconditioning. Journal of Computer and System Sciences, 16(2):124–139, 1978. URL: https://www.sciencedirect.com/science/article/pii/0022000078900417.
\bibitem{ref34} Meena Mahajan. Algebraic Complexity Classes. In Perspectives in Computational Complexity, pages 51–75. Springer, 2014. doi:10.1007/978-3-319-05446-9\_4.
\bibitem{ref35} Carlos Moreira. On asymptotic estimates for arithmetic cost functions. Proceedings of the American Mathematical Society, 125(2):347–353, 1997. URL: https://www.jstor.org/ stable/2161660.
\bibitem{ref36} Ketan D Mulmuley. The GCT program toward the P vs. NP problem. Communications of the ACM, 55(6):98–107, 2012. doi:10.1145/2184319.2184341.
\bibitem{ref37} Danny Nguyen and Igor Pak. Complexity of short generating functions. In Forum of Mathematics, Sigma, volume 6. Cambridge University Press, 2018. arXiv:1702.08660.
\bibitem{ref38} Øystein Ore. Über höhere kongruenzen. Norsk Mat. Forenings Skrifter, 1(7):15, 1922.
\bibitem{ref39} Igor Pak. Complexity problems in enumerative combinatorics. In Proceedings of the Inter- national Congress of Mathematicians – Rio de Janeiro 2018. Vol. IV. Invited lectures, pages 3153–3180. World Sci. Publ., Hackensack, NJ, 2018. doi:10.1142/9789813272880\_0176.
\bibitem{ref40} David A Plaisted. New NP-hard and NP-complete polynomial and integer divisibility problems. Theoretical Computer Science, 31(1-2):125–138, 1984. Preliminary in the 17th Annual Symposium on Foundations of Computer Science (FOCS 1976). URL: https: //www.sciencedirect.com/science/article/pii/0304397584901300.
\bibitem{ref41} Jacob T Schwartz. Fast probabilistic algorithms for verification of polynomial identities. Journal of the ACM (JACM), 27(4):701–717, 1980. doi:10.1145/322217.322225.
\bibitem{ref42} Adi Shamir. Factoring numbers in O (logn) arithmetic steps. Information Processing Let- ters, 8(1):28–31, 1979. URL: https://www.sciencedirect.com/science/article/abs/pii/ 0020019079900875.
\bibitem{ref43} Amir Shpilka and Amir Yehudayoff. Arithmetic Circuits: A survey of recent results and open questions. Foundations and Trends® in Theoretical Computer Science, 5(3–4):207–388, 2010. doi:10.1561/0400000039.
\bibitem{ref44} Michael Shub and Steve Smale. On the intractability of Hilbert’s Nullstellensatz and an algebraic version of “NP $\neq$ P?”. Duke Math. J., 81(1):47–54 (1996), 1995. A celebration of John F. Nash, Jr. doi:10.1215/S0012-7094-95-08105-8.
\bibitem{ref45} Volker Strassen. Vermeidung von divisionen. Journal für die reine und angewandte Mathematik, 264:184–202, 1973.
\bibitem{ref46} L Valiant. Reducibility by algebraic projections in: Logic and algorithmic. In Symposium in honour of Ernst Specker, pages 365–380, 1982.
\bibitem{ref47} Leslie G Valiant. Completeness classes in algebra. In Proceedings of the 11th Annual ACM symposium on Theory of computing, pages 249–261. ACM, 1979. doi:10.1145/800135.804419.
\bibitem{ref48} Joachim Von Zur Gathen and Jürgen Gerhard. Modern computer algebra. Cambridge university press, 2013.
\bibitem{ref49} Wact. Some Accessible Open Problems. Workshop on Algebraic Complexity Theory (WACT 2016).   URL: https://www.cs.tau.ac.il/~shpilka/wact2016/concreteOpenProblems/ openprobs.pdf.
\bibitem{ref50} Klaus W Wagner. The complexity of combinatorial problems with succinct input representation. Acta informatica, 23(3):325–356, 1986. doi:10.1007/BF00289117.
\bibitem{ref51} Richard Zippel. Probabilistic Algorithms for Sparse Polynomials. In Proceedings of the International Symposium on Symbolic and Algebraic Computation, EUROSAM ’79, pages 216–226, 1979. doi:10.1007/3-540-09519-5\_73.
\bibitem{ref52} J.von zur Gathen and V. Strassen. Some polynomials that are hard to compute. Theoretical Computer Science, 11(3):331–335, 1980. URL: http://www.sciencedirect.com/science/ article/pii/0304397580900201.
\end{thebibliography}
\end{document}
